\documentclass[11pt,class=report,crop=false]{standalone}
\usepackage[screen]{../logic}


\begin{document}


%====================================================================
\chapitre{Logique vrai/faux}
%====================================================================

% \insertvideo{}{partie 1.1. Titre}


\objectifs{La logique \og{}Vrai/Faux\fg{}, aussi appelée logique propositionnelle, s'attache à étudier quelle valeur, parmi \og{}Vrai\fg{} ou \og{}Faux\fg{}, peut prendre une formule définie à partir des connecteurs logiques \tand, \tor, \tnot\ldots}


%%%%%%%%%%%%%%%%%%%%%%%%%%%%%%%%%%%%%%%%%%%%%%%%%%%%%%%%%%%%%%%%%%%%%
\section{Vrai/faux, connecteurs}

%--------------------------------------------------------------------
\subsection{Vrai/faux}

Nous définissons deux valeurs \ltrue{}\index{001@$\ltrue$} (Vrai) et \lfalse{}\index{002@$\lfalse$} (Faux), appelées \defi{valeurs booléennes} ou \defi{valeurs de vérité}.
On utilisera aussi la notation \lzero{}\index{003@$\lzero$} (Faux) et \lone{}\index{004@$\lone$} (Vrai).
La notation anglo-saxonne est \ltruetrue{} (\emph{True}) et \lfalsefalse{} (\emph{False}).


Une \defi{variable propositionnelle}\index{variable!propositionnelle} est une variable qui peut prendre soit la valeur \ltrue{} (Vrai), soit la valeur \lfalse{} (Faux).

Par exemple, si $P$ : \og{}$2+2 = 4$\fg{}, alors $P$ est vraie ; en revanche, $Q$ : \og{}$2^2 < 3$\fg{} est fausse.


%--------------------------------------------------------------------
\subsection{$\lnot$ (\tnot), $\land$ (\tand), $\lor$ (\tor)}

À partir d'une ou plusieurs variables propositionnelles et de \defi{connecteurs}, on forme de nouvelles formules.
Soient $P$ et $Q$ deux variables propositionnelles.

\begin{itemize}
	\item \defi{Négation}\index{negation@négation} $\lnot P$.\index{005@$\lnot$}\index{non@\tnot{} logique} \quad $\lnot P$ est vraie si et seulement si $P$ est fausse.
    
	\item \defi{Conjonction}\index{conjonction} $P \land Q$.\index{006@$\land$}\index{et@\tand{} logique} \quad  $P \land Q$ est vraie si et seulement si $P$ et $Q$ sont vraies.
    
	\item \defi{Disjonction}\index{disjonction} $P \lor Q$.\index{007@$\lor$}\index{ou@\tor{} logique} \quad $P \lor Q$ est vraie si et seulement si $P$ est vraie ou $Q$ est vraie.
\end{itemize}



Reprenons $P$ : \og{}$2+2 = 4$\fg{} et $Q$ : \og{}$2^2 < 3$\fg{}.
Alors : 
\begin{itemize}
	\item $\lnot P$ est fausse (car $P$ est vraie),
	\item $\lnot Q$ est vraie (car $Q$ est fausse),	
	\item $P \land Q$ est fausse (car $Q$ n'est pas vraie),
	\item $P \lor Q$ est vraie (car $P$ est vraie).
\end{itemize}

Évidemment, il y a un problème de logique dans nos définitions, car par exemple nous avons défini le connecteur \og{}$\land$ \fg{} (\tand) à l'aide du mot français \og{}et\fg{}. Pour remédier à ce problème, on définit l'action de chacun de ces connecteurs par sa \defi{table de vérité}\index{table de verite@table de vérité} :

\[
\begin{array}{c | c}
P		& \lnot P \\
\hline
&	\\[-2ex]
\ltrue 	& \lfalse \\
\lfalse & \ltrue \\
\end{array}
\qquad\qquad\qquad
\begin{array}{c c | c}
	P		& Q 		& P \land Q \\
	\hline
	&&	\\[-2ex]
	\ltrue 	& \ltrue 	& \ltrue \\
	\ltrue 	& \lfalse 	& \lfalse \\
	\lfalse & \ltrue 	& \lfalse \\	
	\lfalse & \lfalse 	& \lfalse \\		
\end{array}
\qquad\qquad\qquad
\begin{array}{c c | c}
	P		& Q 		& P \lor Q \\
	\hline
	&&	\\[-2ex]
	\ltrue 	& \ltrue 	& \ltrue \\
	\ltrue 	& \lfalse 	& \ltrue \\
	\lfalse & \ltrue 	& \ltrue \\	
	\lfalse & \lfalse 	& \lfalse \\		
\end{array}
\]


Bien sûr, à partir de ces connecteurs on peut composer des formules plus compliquées :
\[
\lnot(P \lor Q)\quad , \quad 
P \land (\lnot P)\quad , \quad 
P \land (Q \lor R)\quad , \quad 
(P \land Q) \lor R\quad , \ \ldots
\]
Les parenthèses sont importantes.


%--------------------------------------------------------------------
\subsection{$\limplies$ (implication)}

L'\defi{implication}\index{implication} $P \limplies Q$\index{008@$\limplies$} est définie par :
\mybox{
$P \limplies Q = (\lnot P) \lor Q$
}
 Comme le connecteur $\lnot$ est prioritaire sur les autres, on écrit simplement : $\lnot P \lor Q$.

On lit \og{}si $P$ alors $Q$\fg{} ou bien \og{}$P$ implique $Q$\fg{}.
Sa table de vérité est :
\[
\begin{array}{c c | c}
	P		& Q 		& P \limplies Q \\
	\hline
	&&	\\[-2ex]
	\ltrue 	& \ltrue 	& \ltrue \\
	\ltrue 	& \lfalse 	& \lfalse \\
	\lfalse & \ltrue 	& \ltrue \\	
	\lfalse & \lfalse 	& \ltrue \\		
\end{array}
\]

L'usage mathématique classique est de noter ce connecteur $P \implies Q$. Mais au fil du temps son usage a été dévoyé pour signifier \og{}donc\fg{}, ce qu'il n'est pas ! On y reviendra dans la Section \ref{sec:connecteursbis}.

Ainsi \og{}$P \limplies Q$\fg{} signifie \og{}si $P$ est vraie alors $Q$ est vraie\fg{}.
Voici quelques exemples :
\begin{itemize}	
	\item \og{}il pleut\fg{} $\limplies$ \og{}le sol est mouillé\fg{} est une implication vraie. Dans ce livre, on évitera les exemples en français pour limiter les formulations ambiguës et les discussions interminables du genre \og{}pourquoi le sol de ma cuisine serait mouillé s'il pleut dehors\fg{}.
		
	\item \og{}2+2=4\fg{} $\limplies$ \og{}$3 \times 3 = 7$\fg{} est une implication fausse.
	
	\item Soit $n$ un entier. \og{}$n$ pair\fg{} $\limplies$ \og{}$n^2$ pair\fg{} est vraie.	
\end{itemize}

\bigskip

Passons aux aspects surprenants de l'implication.
\begin{itemize}	
	\item \og{}2+2=4\fg{} $\limplies$ \og{}un entier est pair ou impair\fg{} est une implication vraie (car si $P$ et $Q$ sont vraies alors $P \limplies Q$ est vraie), même s'il n'y a aucun lien de causalité entre $P$ et $Q$.
	
	\item \og{}$3 \times 3 = 7$\fg{} $\limplies$ \og{}$2+2=4$\fg{} 
	et \og{}$3 \times 3 = 7$\fg{} $\limplies$ \og{}$2+2=5$\fg{}
	sont aussi des implications vraies.
	En effet, dès que $P$ est fausse, alors $P \limplies Q$ est vraie. C'est bien sûr clair par la définition de l'implication via $\lnot P \lor Q$. Il faut accepter cette conséquence de la définition. Nous y reviendrons aussi dans la Section \ref{sec:connecteursbis}.
\end{itemize}


%--------------------------------------------------------------------
\subsection{$\lequiv$ (équivalence)}

L'\defi{équivalence}\index{equivalence@équivalence} $P \lequiv Q$\index{009@$\lequiv$} est définie par la double implication $P \limplies Q$ et $Q \limplies P$ :
\mybox{
	$P \lequiv Q =  (P \limplies Q) \land (Q \limplies P)$
}
On dit \og{}$P$ équivaut à $Q$\fg{} ou bien \og{}$P$ si et seulement si $Q$\fg{}.

Autrement dit, $P \lequiv Q$ est vraie lorsque $P$ et $Q$ sont vraies en même temps ou fausses en même temps.

\[
\begin{array}{c c | c}
	P		& Q 		& P \lequiv Q \\
	\hline
	&&	\\[-2ex]
	\ltrue 	& \ltrue 	& \ltrue \\
	\ltrue 	& \lfalse 	& \lfalse \\
	\lfalse & \ltrue 	& \lfalse \\	
	\lfalse & \lfalse 	& \ltrue \\		
\end{array}
\]

Voici des exemples d'équivalences vraies :
\begin{itemize}
	\item $\lnot \lnot P \ \lequiv \ P$ : la négation d'une négation revient à la proposition initiale.
	
	\item $P \limplies Q \ \lequiv \ (\lnot P) \lor Q$ : par définition de l'implication.
\end{itemize}

En revanche, l'équivalence suivante est \emph{fausse} :
\[
P \land Q \ \lequiv \ P \lor Q
\]
Plus précisément, cette équivalence n'est pas vraie pour toutes les valeurs attribuées à $P$ et $Q$.
En effet, l'implication $P \land Q \limplies P \lor Q$ est toujours vraie, mais l'implication $P \lor Q \limplies P \land Q$ est fausse lorsque $P$ est vraie et $Q$ est fausse par exemple.


%--------------------------------------------------------------------
\subsection{Résumé}

Il y a deux valeurs booléennes \ltrue/\lfalse{} (Vrai/Faux, ou \ltruetrue/\lfalsefalse{} \emph{True}/\emph{False}, ou \lone/\lzero).
Nous avons défini les cinq connecteurs logiques de base.

\begin{center}
\begin{tabular}{| c | c | c |}
    \hline 
	\textbf{Symbole}		& \textbf{Nom} 		& \textbf{Définition} \\
	\hline
	&&	\\[-2ex]
	$\lnot$ 	& \tnot			& $\lnot P$ est vrai ssi $P$ est faux \\
	$\land$ 	& \tand			& $P \land Q$ est vrai ssi $P$ et $Q$ sont vrais \\	
	$\lor$ 		& \tor			& $P \lor Q$ est vrai ssi $P$ ou $Q$ est vrai \\	
	$\limplies$ 	& implication	& $\lnot P \lor Q$ \\	
	$\lequiv$ 	& équivalence	& $(P \limplies Q) \land (Q \limplies P)$ \\
    \hline				
\end{tabular}
\end{center}

Le connecteur $\lnot$ est un connecteur\index{connecteur} \defi{unaire}, car il s'applique à un seul paramètre (comme une fonction d'une variable) alors que les autres sont des connecteurs \defi{binaires} (par exemple on pourrait écrire $P \land Q$ sous la forme $\land(P,Q)$ comme une fonction de deux variables).



%%%%%%%%%%%%%%%%%%%%%%%%%%%%%%%%%%%%%%%%%%%%%%%%%%%%%%%%%%%%%%%%%%%%%
\section{Formules}

%--------------------------------------------------------------------
\subsection{Objectifs}

Nous disposons de \defi{variables propositionnelles}\index{variable!propositionnelle} $P$, $Q$, $R$\ldots{} chacune pouvant prendre la valeur Vrai ou Faux. À l'aide des connecteurs, on compose une \emph{formule} $\mathcal{P}$, par exemple :
\[
\mathcal{P} :  \quad P \land \big( Q \lor (\lnot R) \big)
\]

Voici les questions qu'on se pose.
Si on remplace $P$, $Q$, $R$\ldots{} par des valeurs \ltrue/\lfalse, est-ce que la formule $\mathcal{P}$ est vraie ou fausse ? On peut même être plus précis : 
\begin{itemize}
	\item Est-ce qu'il existe un choix de valeurs \ltrue/\lfalse{} qui rende l'évaluation de $\mathcal{P}$ vraie ?
	\item Pour tous les choix possibles de valeurs \ltrue/\lfalse, l'évaluation de $\mathcal{P}$ est-elle vraie ?
\end{itemize}


Par exemple, pour la proposition $\mathcal{P}$ : $P \land \big( Q \lor (\lnot R) \big)$, si on attribue la valeur \ltrue{} à $P$, \lfalse{} à $Q$ et \lfalse{} à $R$, alors $\mathcal{P}$ s'évalue en $\ltrue \land ( \lfalse \lor (\lnot \lfalse))$ qui vaut \ltrue.
En revanche, si on remplace $(P,Q,R)$ par $(\ltrue,\lfalse,\ltrue)$, alors $\mathcal{P}$ s'évalue en \lfalse.

Bien sûr, on pourrait se poser les mêmes questions pour savoir quand $\mathcal{P}$ s'évalue en faux, mais cela revient à décider quand la formule $\lnot \mathcal{P}$ s'évalue en vrai.

Si on dispose de deux formules $\mathcal{P}$ et $\mathcal{Q}$, alors on peut composer des formules plus complexes $\mathcal{P} \land \mathcal{Q}$, $\mathcal{P} \lor \mathcal{Q}$, etc.

Lorsque l'on dispose de deux formules $\mathcal{P}$ et $\mathcal{Q}$, l'objectif le plus important sera de comprendre le lien entre $\mathcal{P}$ et $\mathcal{Q}$. Par exemple, est-ce que $\mathcal{P} \limplies \mathcal{Q}$ est vraie ? Est-ce que $\mathcal{P} \lequiv \mathcal{Q}$ est vraie ?


%--------------------------------------------------------------------
\subsection{Définition d'une formule}

On considère un \defi{alphabet}\index{alphabet} formé par :
\begin{itemize}
	\item un ensemble $\text{Var} = \{P, Q, R,\ldots\}$ de variables propositionnelles,
	
	\item un ensemble $\text{Conn} = \{ \lnot, \land, \lor, \limplies, \lequiv \}$ de connecteurs logiques,
	
	\item de symboles de parenthèses $($ et $)$\,.
\end{itemize}

Une \defi{formule}\index{formule} est un mot fini, formé à partir des caractères de cet alphabet ; ce mot doit être correctement parenthésé et cohérent avec l'usage des connecteurs.

\begin{definition}
Plus précisément, l'ensemble $\text{Form}$ de toutes les \defi{formules} est le plus petit ensemble de mots, contenant les variables propositionnelles, tel que :
\begin{itemize}
	\item si $\mathcal{P}$ est une formule, alors $\lnot \mathcal{P}$ aussi,
	\item si $\mathcal{P}$ et $\mathcal{Q}$ sont des formules, alors 
	$(\mathcal{P} \land \mathcal{Q})$, 
	$(\mathcal{P} \lor \mathcal{Q})$, 
	$(\mathcal{P} \limplies \mathcal{Q})$, 
	$(\mathcal{P} \lequiv \mathcal{Q})$ sont aussi des formules.
\end{itemize}
\end{definition}

Remarque : on notera $P, Q, R\ldots$ (caractère droit) pour les variables propositionnelles et 
$\mathcal{P}, \mathcal{Q}, \mathcal{R}\ldots$ (caractère calligraphique) pour les formules.

Voici des exemples de formules :
\[
P \quad, \quad 
\lnot P \quad, \quad 
(Q \limplies P) \quad, \quad 
((Q \limplies P) \lequiv Q) \quad, \quad 
(\lnot R \limplies (P \lor Q))
\]

Par abus de notation :
\begin{itemize}
	\item On se permettra d'enlever les parenthèses quand elles sont totalement à l'extérieur, par exemple on écrira $Q \limplies P$  au lieu de $(Q \limplies P)$ ou lorsqu'il n'y a pas d'ambiguïté, par exemple $P \lor Q \lor R$.
	
	\item À l'inverse, on rajoutera parfois des parenthèses superflues. Par exemple, on pourra écrire $(\lnot P) \land Q$ au lieu de $\lnot P \land Q$ afin d'éviter la confusion avec $\lnot (P \land Q)$.
\end{itemize}

%--------------------------------------------------------------------
\subsection{Table de vérité, évaluation}

Pour une formule $\mathcal{P}$ qui dépend de deux variables propositionnelles $P$ et $Q$, on peut déterminer quand la formule prend la valeur Vrai ou Faux selon les valeurs attribuées à $P$ et $Q$. Cela se résume sous la forme d'une \defi{table de vérité}\index{table de verite@table de vérité} de la formule $\mathcal{P}$ :

\[
\begin{array}{c c | c}
	P		& Q 		& \mathcal{P} \\
	\hline
	&&	\\[-2ex]
	\ltrue 	& \ltrue 	& \ltrue/\lfalse\ ? \\
	\ltrue 	& \lfalse 	& \ltrue/\lfalse\ ? \\
	\lfalse & \ltrue 	& \ltrue/\lfalse\ ? \\	
	\lfalse & \lfalse 	& \ltrue/\lfalse\ ? \\		
\end{array}
\]

Plus généralement, si la formule $\mathcal{P}$ dépend des variables $P_1, P_2,\ldots,P_n$,
une \defi{distribution de valeurs}\index{distribution de valeurs} est un choix \ltrue/\lfalse{} pour la variable $P_1$, un choix \ltrue/\lfalse{} pour la variable $P_2$,\ldots{}
On obtient alors une \defi{évaluation}\index{evaluation@évaluation} \ltrue/\lfalse{} pour la formule $\mathcal{P}$.

Dans la pratique on dresse la table de vérité de la formule $\mathcal{P}$.

\begin{exemple}
Soit :
\[
\mathcal{P} : \qquad (P \land Q) \lor (\lnot P)
\]
On dresse la table de vérité de $\mathcal{P}$ à l'aide des tables intermédiaires de $P \land Q$ et $\lnot P$.

\[
\begin{array}{c c | c | c | c}
	P		& Q 		& P \land Q 	& \lnot P 	&  \mathcal{P} \\
	\hline
	&&&&	\\[-2ex]
	\ltrue 	& \ltrue 	& \ltrue		& \lfalse 	& \ltrue \\
	\ltrue 	& \lfalse 	& \lfalse		& \lfalse	& \lfalse \\
	\lfalse & \ltrue 	& \lfalse		& \ltrue	& \ltrue \\	
	\lfalse & \lfalse 	& \lfalse 		& \ltrue	& \ltrue \\		
\end{array}
\]

\end{exemple}

Dans les chapitres \og{}Quels connecteurs sont utiles ?\fg{} et \og{}Théorème de compacité\fg{}, on présentera une distribution de valeurs comme une fonction $v : \{ \ltrue, \lfalse \}^n \to \{ \ltrue, \lfalse \}$ qui s'étend sur l'ensemble des formules afin d'obtenir une évaluation $v(\mathcal{P})$.


%--------------------------------------------------------------------
\subsection{Équivalence logique}

\begin{definition}
	Deux formules $\mathcal{P}$ et $\mathcal{Q}$ sont \defi{logiquement équivalentes}\index{equivalence logique@équivalence logique}\index{010@$\equiv$} si elles ont la même table de vérité. On note alors :
	\[ \mathcal{P} \equiv \mathcal{Q} \]
\end{definition}
Cela signifie que les formules $\mathcal{P}$ et $\mathcal{Q}$ s'évaluent à Vrai en même temps et à Faux en même temps, quel que soit le choix des valeurs de vérité.

Il existe un lien direct avec le connecteur d'équivalence $\lequiv$. En effet, $\mathcal{P} \equiv \mathcal{Q}$ signifie que $(\mathcal{P} \lequiv \mathcal{Q})$ est une tautologie (c'est-à-dire toujours évaluée à Vrai), voir la Section \ref{sec:formulesbis}.

\begin{exemple}
	\sauteligne
	\begin{enumerate}
		\item $\lnot \lnot P \equiv P$ (double négation)
		\item $\lnot (P \lor Q) \equiv (\lnot P) \land (\lnot Q)$ (loi de De Morgan)
		\item $(P \land Q) \limplies R \equiv \big( \lnot (P \land Q)\big) \lor R \equiv \big( (\lnot P) \lor (\lnot Q) \big) \lor R$
	\end{enumerate}
\end{exemple}


%--------------------------------------------------------------------
\subsection{Équivalences logiques usuelles}


\[
\begin{array}{l l l}
	\text{idempotence} 
	& \mathcal{P}\land\mathcal{P} \equiv \mathcal{P}
	& \mathcal{P}\lor\mathcal{P} \equiv \mathcal{P}
	\\[0.6ex]
	
	\text{commutativité}
	& \mathcal{P}\land\mathcal{Q} \equiv \mathcal{Q}\land\mathcal{P}
	& \mathcal{P}\lor\mathcal{Q} \equiv \mathcal{Q}\lor\mathcal{P}
	\\[0.6ex]
	
	\text{associativité}
	& (\mathcal{P}\land\mathcal{Q})\land\mathcal{R} \equiv \mathcal{P}\land(\mathcal{Q}\land\mathcal{R})
	& (\mathcal{P}\lor\mathcal{Q})\lor\mathcal{R} \equiv \mathcal{P}\lor(\mathcal{Q}\lor\mathcal{R})
	\\[0.6ex]
	
	\text{distributivité}
	& \mathcal{P}\land(\mathcal{Q}\lor\mathcal{R}) \equiv (\mathcal{P}\land\mathcal{Q}) \lor (\mathcal{P}\land\mathcal{R})
	& \mathcal{P}\lor(\mathcal{Q}\land\mathcal{R}) \equiv (\mathcal{P}\lor\mathcal{Q}) \land (\mathcal{P}\lor\mathcal{R})
	\\[0.6ex]
	
	\text{absorption}
	& \mathcal{P}\land(\mathcal{P}\lor\mathcal{Q}) \equiv \mathcal{P}
	& \mathcal{P}\lor(\mathcal{P}\land\mathcal{Q}) \equiv \mathcal{P}
	\\[0.6ex]
	
	\text{lois de De Morgan}\index{lois de De Morgan}
	& \lnot(\mathcal{P}\land\mathcal{Q}) \equiv (\lnot\mathcal{P})\lor(\lnot\mathcal{Q})
	& \lnot(\mathcal{P}\lor\mathcal{Q}) \equiv (\lnot\mathcal{P})\land(\lnot\mathcal{Q})
	\\[0.6ex]
	
	\text{contraposée}\index{contraposition}
	& \mathcal{P}\limplies\mathcal{Q} \equiv (\lnot\mathcal{Q})\limplies(\lnot\mathcal{P})
	& \\
\end{array}
\]

\medskip

Ce sont les plus classiques. Il en existe beaucoup d'autres !



%%%%%%%%%%%%%%%%%%%%%%%%%%%%%%%%%%%%%%%%%%%%%%%%%%%%%%%%%%%%%%%%%%%%%
\section{Connecteurs (suite)}
\label{sec:connecteursbis}
\index{connecteur}
Nous apportons des précisions et des commentaires sur nos connecteurs favoris.

%--------------------------------------------------------------------
\subsection{Vrai/Faux, négation}
\index{negation@négation}
\index{non@\tnot{} logique}

Il y a exactement deux valeurs booléennes $\{\ltrue, \lfalse\}$ (ou $\{\lzero, \lone\}$), ni plus, ni moins ! Ce qui fait que si une variable $P$ ou une formule $\mathcal{P}$ n'est pas évaluée à Faux, c'est qu'elle vaut Vrai. Et réciproquement. Il n'y a donc pas de réponse \og{}peut-être\fg{}.
Cela a les conséquences suivantes :
\begin{itemize}
	\item \defi{Principe du tiers exclu}\index{tiers exclu} :
\mybox{$\mathcal{P} \lor \lnot \mathcal{P} \quad \text{ est toujours Vrai }$}

L'expression \og{}tiers exclu\fg{} signifie qu'il n'y a pas d'autre choix que \ltrue/\lfalse.

	\item Principe de non-contradiction :
\[ \mathcal{P} \land \lnot \mathcal{P} \quad \text{ est toujours Faux } \]

	\item Loi de la double négation :
\[ \lnot\lnot \mathcal{P} \equiv \mathcal{P}\]
	
\end{itemize}

On verra un peu plus tard dans la Section \ref{sec:formulesbis} une nouvelle écriture de ce qui est vrai/faux :\index{011@$\top$}
\[ \top \quad : \text{ formule qui s'évalue toujours à Vrai }\]
Ce symbole désigne une formule qui s'évalue toujours à \ltrue{} (\ltruetrue{} en anglais, d'où le symbole). Il se lit \og{}toujours Vrai\fg{} ou \og{}tautologie\fg{} ou \emph{top}.

Sa contrepartie est le symbole renversé :\index{012@$\bot$}
\[ \bot \quad : \text{ formule qui s'évalue toujours à Faux }\]

Ainsi on peut écrire :
\[ \mathcal{P} \lor \lnot \mathcal{P} \equiv \top \qquad \mathcal{P} \land \lnot \mathcal{P} \equiv \bot\]


Quelle est la différence entre $\ltrue$ et $\top$ ?
$\ltrue$ est une valeur, alors que $\top$ est une formule qui s'évalue toujours à $\ltrue$.
Si on identifie plutôt les valeurs booléennes à $\lzero/\lone$, alors $\lone$ est une constante, tandis que $\top$ est la fonction constante égale à $\lone$.

%--------------------------------------------------------------------
\subsection{$\land$ (\tand), $\lor$ (\tor)}
\index{conjonction}\index{et@\tand{} logique}
\index{disjonction}\index{ou@\tor{} logique}

Dans notre liste d'équivalences logiques, nous avions :
\[ \mathcal{P} \land \mathcal{Q} \equiv \mathcal{Q} \land \mathcal{P}\qquad \mathcal{P} \lor \mathcal{Q} \equiv \mathcal{Q} \lor \mathcal{P} \qquad\qquad \text{commutativité} \]
et aussi :
\[ (\mathcal{P} \land \mathcal{Q}) \land \mathcal{R} \equiv \mathcal{P} \land (\mathcal{Q} \land \mathcal{R}) 
\qquad 
(\mathcal{P} \lor \mathcal{Q}) \lor \mathcal{R} \equiv \mathcal{P} \lor (\mathcal{Q} \lor \mathcal{R}) 
\qquad\qquad \text{associativité} \]
Par conséquent, on s'autorise à écrire $\mathcal{P} \land \mathcal{Q} \land \mathcal{R}$ et $\mathcal{P} \lor \mathcal{Q} \lor \mathcal{R}$, sans parenthèses.
Enfin :
\[ \mathcal{P} \land (\mathcal{Q} \lor \mathcal{R}) \equiv (\mathcal{P} \land \mathcal{Q}) \lor (\mathcal{P} \land \mathcal{R}) \qquad \mathcal{P} \lor (\mathcal{Q} \land \mathcal{R}) \equiv (\mathcal{P} \lor \mathcal{Q}) \land (\mathcal{P} \lor \mathcal{R}) \qquad\qquad \text{distributivité} \]

\medskip

Le connecteur $\land$ (\tand) est analogue à la multiplication $\times$, tandis que $\lor$ (\tor) est analogue à l'addition $+$ :
$a \times b = b \times a$ (commutativité), $(a \times b) \times c = a \times (b \times c)$ (associativité), $a \times (b+c) = a \times b + a \times c$ (distributivité).

Les lois de De Morgan se résument par l'incantation \og{}la négation d'un \tand{} est un \tor\fg{} et \og{}la négation d'un \tor{} est un \tand\fg{} :
\[ \lnot(\mathcal{P} \land \mathcal{Q}) \equiv \lnot\mathcal{P} \lor \lnot\mathcal{Q}
\qquad\qquad \lnot(\mathcal{P} \lor \mathcal{Q}) \equiv \lnot\mathcal{P} \land \lnot\mathcal{Q}
 \]

%--------------------------------------------------------------------
\subsection{$\limplies$ (implication)}
\index{implication}

Par définition :
\[
\mathcal{P} \limplies \mathcal{Q} \equiv \lnot \mathcal{P} \lor \mathcal{Q}
\]
et sa table de vérité est :
\[
\begin{array}{c c | c}
	\mathcal{P}		& \mathcal{Q} 		& \mathcal{P} \limplies \mathcal{Q} \\
	\hline
	&&	\\[-2ex]
	\ltrue 	& \ltrue 	& \ltrue \\
	\ltrue 	& \lfalse 	& \lfalse \\
	\lfalse & \ltrue 	& \ltrue \\	
	\lfalse & \lfalse 	& \ltrue \\		
\end{array}
\]

Il y a plusieurs points à intégrer concernant les implications.

\medskip

\textbf{Pas de causalité.} 

Il faut bien comprendre que $P \limplies Q$ (ici avec des variables propositionnelles) ne signale pas un lien de causalité mais plutôt une constatation : est-ce que si $P$ est vraie alors $Q$ est aussi vraie ?


Voici un exemple qui devrait faire hurler tous vos profs de maths.
\begin{exemple}
 Soit $n$ un entier :
\[
n \text { est impair} \limplies n \text{ est premier}
\]
Tous vos profs vous diront que cette implication est fausse, mais ce n'est pas vrai !
La confusion vient du fait qu'on a l'habitude de comprendre cette phrase par \og{}tout nombre impair est aussi un nombre premier\fg{}, ce qui est bien sûr faux.
Mais ici $n$ est fixé et sa valeur n'était pas précisée et donc l'implication peut être vraie ou fausse selon la valeur de $n$.

\begin{itemize}
	\item si $n=3$, alors \og{}$n \text { est impair} \limplies n \text{ est premier}$\fg{} est vraie, car $P$ : \og{}$3$ est impair\fg{} et $Q$ : \og{}$3$ est premier\fg{} sont vraies.
	
	\item si $n=15$, alors \og{}$n \text { est impair} \limplies n \text{ est premier}$\fg{} est fausse, car $P$ : \og{}$15$ est impair\fg{} est vraie alors que $Q$ : \og{}$15$ est premier\fg{} est fausse.
	
	\item si $n=4$, alors \og{}$n \text { est impair} \limplies n \text{ est premier}$\fg{} est vraie ! En effet $P$ : \og{}$4$ est impair\fg{} est fausse et donc $P \limplies Q$ est vraie. 
\end{itemize}


\end{exemple}


\medskip

\textbf{Lorsque $P$ est faux alors $P \limplies Q$ est vrai !}

C'est bien sûr clair d'après la définition $(\lnot P) \lor Q$.
Essayons de justifier pourquoi c'est naturel. 
Posons-nous la question : \og{}quand est-ce qu'on veut que $P$ implique $Q$ soit faux ?\fg{}. La réponse est : \og{}si $P$ est vrai et $Q$ est faux\fg{} (c'est comme avoir trouvé un contre-exemple). Dans tous les autres cas, $P \limplies Q$ est vrai, donc y compris lorsque $P$ est faux.

Cela se retrouve dans le langage courant :
\mycenterline{
	\og{}Si tu as raison, alors moi je suis le pape.\fg
}
Sous-entendu, \og{}Tu as tort.\fg{}
Dans le même genre :
\mycenterline{
	\og{}Je le ferai quand les poules auront des dents.\fg
}
Sous-entendu, \og{}Je ne le ferai pas.\fg{} C'est l'implication \og{}poules avec dents $\limplies$ action\fg{} qui est toujours vraie avec ou sans action.


\medskip

\textbf{À propos des démonstrations.}

De $P \limplies Q$, on ne peut pas déduire que $Q$ est vrai. Ce que l'on verra dans la suite s'appelle le \emph{modus ponens} et est la base de nos raisonnements :
\mycenterline{
	Si $P$ est vrai et si $P\limplies Q$ est vrai, alors $Q$ est vrai.
}

C'est l'erreur que l'on fait à force d'habitude quand on écrit $P \implies Q$ : on suppose implicitement que $P$ est vrai et on en déduit que $Q$ est vrai.

Enfin $\mathcal{P} \limplies \mathcal{Q} \limplies \mathcal{R}$ ne dénote pas une chaîne de conséquences, mais c'est par définition la formule :
\[ \mathcal{P} \limplies (\mathcal{Q} \limplies \mathcal{R}) \]



%--------------------------------------------------------------------
\subsection{$\lequiv$ (équivalence)}
\index{equivalence@équivalence}

Pour qu'une équivalence soit vraie, il faut que $\mathcal{P}$ et $\mathcal{Q}$ soient vraies en même temps, et fausses en même temps.
Par définition :
\[ \mathcal{P} \lequiv \mathcal{Q} \equiv (\mathcal{P} \limplies \mathcal{Q}) \land (\mathcal{Q} \limplies \mathcal{P}) \]

On peut formuler les mêmes avertissements que pour l'implication : quand on écrit $\mathcal{P} \lequiv \mathcal{Q}$ cela ne signifie pas que $\mathcal{P}$ et $\mathcal{Q}$ sont toujours évaluées à vrai, mais juste qu'elles s'évaluent à vrai en même temps.
 

%--------------------------------------------------------------------
\subsection{D'autres connecteurs}
\index{connecteur}

Un \defi{connecteur à $n$ places} est une fonction $f : \{\ltrue, \lfalse\}^n \to \{\ltrue, \lfalse\}$, ou, si vous préférez, une fonction $f : \{\lzero, \lone\}^n \to \{\lzero, \lone\}$.

\begin{itemize}
	\item $n=1$. $\lnot$ est un connecteur unaire (c'est une fonction à une variable). On peut aussi avoir besoin de l'identité $\mathrm{id}$ qui envoie $\ltrue$ sur $\ltrue$ et $\lfalse$ sur $\lfalse$.
	
	\item $n=2$. $\land$ est un connecteur binaire (comme une fonction de deux variables). Sauf qu'on écrit $P \land Q$ au lieu de $\land(P, Q)$.
	
	\item $n=2$. Le \defi{ou exclusif}\index{015@$\oplus$}\index{ou exclusif@\tconnector{xor} ou exclusif} $\oplus$/\tconnector{xor} est un autre connecteur binaire :
	\mycenterline{$P \oplus Q$ est vrai si exactement l'une des variables est vraie.}
	
	\item $n=2$. Le connecteur $\uparrow$/\tconnector{nand},\index{016@$\uparrow$}\index{nand@\tconnector{nand}} aussi appelé symbole de Sheffer, est défini comme la négation du \tand :
	\[ P \uparrow Q = \lnot (P \land Q) \]
	On listera tous les connecteurs binaires dans le chapitre \og{}Quels connecteurs sont utiles ?\fg.
	
	\item Pour $n$ quelconque, on peut définir un connecteur à $n$ places, en donnant la table de vérité en fonction des valeurs affectées aux variables $P_1,P_2,\ldots, P_n$. Cette table a $2^n$ lignes.
	On verra, toujours dans le chapitre \og{}Quels connecteurs sont utiles ?\fg{}, pourquoi on n'a pas vraiment besoin de ces connecteurs.
\end{itemize}



%%%%%%%%%%%%%%%%%%%%%%%%%%%%%%%%%%%%%%%%%%%%%%%%%%%%%%%%%%%%%%%%%%%%%
\section{Formules (suite)}
\label{sec:formulesbis}

%--------------------------------------------------------------------
\subsection{Formules satisfaisables}

Soient $P_1, P_2,\ldots, P_n$ des variables propositionnelles.
Une distribution de valeurs $v$ associe à chaque variable $P_i$ une valeur de vérité $v(P_i)$ qui est \ltrue{} ou \lfalse{} (ou une valeur \lzero/\lone).
On peut alors associer une valeur \ltrue/\lfalse, notée $v(\mathcal{P})$, à n'importe quelle formule $\mathcal{P}$.

\begin{exemple}
On considère trois variables $P, Q, R$ et la formule :
\[ \mathcal{P} : \quad (P \land \lnot Q) \lor (Q \land R) \]
Considérons la distribution de valeurs $(\ltrue, \lfalse, \lfalse)$ (c'est-à-dire $v(P) = \ltrue$, $v(Q)=\lfalse$, $v(R)=\lfalse$), alors, lorsque l'on substitue ces valeurs dans $\mathcal{P}$, on obtient :
\[ v(\mathcal{P}) \ = \  (\ltrue \land \lnot \lfalse) \lor (\lfalse \land \lfalse) \]
et donc $v(\mathcal{P}) = \ltrue$.
Par contre, pour la distribution $(\ltrue,\ltrue,\lfalse)$ alors $\mathcal{P}$ s'évalue à \lfalse.
\end{exemple}

\begin{definition}
	\sauteligne
	\index{formule!satisfaction}
	\begin{itemize}
		\item Une formule $\mathcal{P}$ est \defi{satisfaite} pour la distribution de valeurs $v$ si $\mathcal{P}$ s'évalue à Vrai pour cette distribution. 
		
		\item Une formule $\mathcal{P}$ est \defi{satisfaisable}, s'il existe une distribution de valeurs pour laquelle $\mathcal{P}$ s'évalue à Vrai.

	\end{itemize}	
\end{definition}


La formule $\mathcal{P} :  (P \land \lnot Q) \lor (Q \land R)$ de l'exemple précédent est satisfaisable, car on avait trouvé une distribution de valeurs qui rend $\mathcal{P}$ vraie. 


%--------------------------------------------------------------------
\subsection{Tautologie}

\begin{definition}
	\sauteligne
	\begin{itemize}
		\item Une \defi{tautologie}\index{tautologie} est une formule $\mathcal{P}$ qui s'évalue toujours à Vrai, quelle que soit la distribution des valeurs de vérité.
		
		\item Une \defi{contradiction} est une formule $\mathcal{P}$ qui s'évalue toujours à Faux, quelle que soit la distribution des valeurs de vérité.
	\end{itemize}
\end{definition}

Notation :
\mybox{$
\text{tautologie} : \ \top \qquad\qquad \quad 
\text{contradiction} : \  \bot
$}

Quelques remarques : 
\begin{itemize}
	\item Ainsi une tautologie est une formule toujours vraie (plus précisément qui s'évalue toujours à la valeur Vrai). 
	Le terme vient du grec \emph{tautos} (le même) et \emph{logos} (parole) : c'est toujours la même valeur (Vrai).
	\item Une contradiction est une formule toujours fausse. 
	\item Cependant, la plupart des formules ne sont ni des tautologies, ni des contradictions : elles s'évaluent à des valeurs vraies et fausses selon le choix de la distribution de valeurs. 
	\item Une formule $\mathcal{P}$ est une contradiction si et seulement si $\lnot\mathcal{P}$ est une tautologie. C'est pourquoi on s'intéressera essentiellement aux tautologies.
\end{itemize}

\begin{exemple}
	\sauteligne
	\begin{enumerate}
		\item $P \lor \lnot P$ est une tautologie (si $P$ est Vrai, $P \lor \lnot P$ est Vrai ; si $P$ est Faux, $P \lor \lnot P$ est aussi Vrai).
		
		\item $P \land \lnot P$ est une contradiction (c'est une formule toujours fausse, quelle que soit la valeur attribuée à $P$).
		
		\item $(P \lor Q \lor R) \lor (\lnot P \land \lnot Q \land \lnot R)$ est une tautologie.
		
		\item $(P \land Q) \limplies Q$ est une tautologie.
		
		\item $\lnot(P \land Q) \lequiv (\lnot P) \lor (\lnot Q)$ est une tautologie.
		
		\item $\big( (P \limplies Q) \land (Q \limplies R) \big) \limplies (P \limplies R)$ est une tautologie.
	\end{enumerate}
\end{exemple}	

Voici des exemples d'utilisation des symboles $\top/\bot$ :
\[
\mathcal{P} \land \top \equiv \mathcal{P} \qquad
\mathcal{P} \lor \top \equiv \top \qquad
\mathcal{P} \land \bot \equiv \bot \qquad
\mathcal{P} \lor \bot \equiv \mathcal{P}
\]


Nous pouvons maintenant préciser le lien entre l'équivalence logique ($\equiv$) et l'équivalence ($\lequiv$) :
\[
\mathcal{P} \equiv \mathcal{Q} \quad \text{ ssi } \quad
\mathcal{P} \lequiv \mathcal{Q} \text{ est une tautologie}
\]
Si c'est le cas, c'est-à-dire si $\mathcal{P} \equiv \mathcal{Q}$ (ou bien si $\mathcal{P} \lequiv \mathcal{Q}$ est toujours vraie), alors $\mathcal{P}$ et $\mathcal{Q}$ s'évaluent à vrai en même temps et à faux en même temps, quelle que soit l'évaluation.

\medskip
Dans les sections suivantes, nous allons découvrir trois autres visions des formules et de leurs évaluations.

%--------------------------------------------------------------------
\subsection{Portes logiques et circuits}
\index{circuit logique}

Pour se rapprocher de ce qu'il se passe au cœur d'un ordinateur on peut considérer les connecteurs comme des portes logiques. Ici les valeurs de vérité sont $\{\lzero, \lone\}$.


\textbf{Porte \tnot/\tconnector{not}.}

Elle a une seule entrée et change \lzero{} en \lone, et \lone{} en \lzero.
\myfigure{1}{
	\tikzinput{portes_01}
}

\textbf{Porte \tand/\tconnector{and}.}

\myfigure{0.8}{
	\tikzinput{portes_02}
}

On résume l'action de la porte \tconnector{and} par le tableau suivant :

\begin{center}
	\begin{minipage}{0.45\textwidth}
		\[			
			\begin{array}{cc|c}
				\hline
				P & Q & \text{Sortie} \\ \hline
				&&	\\[-2ex]
				\lzero & \lzero & \lzero \\
				\lzero & \lone & \lzero \\
				\lone & \lzero & \lzero \\
				\lone & \lone & \lone \\ 
			\end{array}
	\]
	\end{minipage}
	\begin{minipage}{0.45\textwidth}
		\myfigure{1}{
			\tikzinput{portes_02b}
		}
	\end{minipage}
\end{center}

\textbf{Porte \tor/\tconnector{or}.}

\myfigure{0.8}{
	\tikzinput{portes_03}
}


\textbf{Porte \tconnector{nand}.}
	
Elle peut se réaliser comme la porte \tconnector{not}(\tconnector{and}), c'est-à-dire la porte \tconnector{and} suivie de la porte \tconnector{not}.
\myfigure{1}{
	\tikzinput{portes_05a}
}
Mais nous préférons la voir comme une porte à part entière. 

\myfigure{1}{
	\tikzinput{portes_05b}
}


\begin{exemple}
Un \emph{multiplexeur} est un circuit qui aiguille un signal entre deux entrées $P$ et $Q$ à l'aide d'un sélecteur $I$.
Si $I$ vaut $\lone$ alors la sortie du circuit est la valeur de $P$.
Si $I$ vaut $\lzero$ alors la sortie est la valeur de $Q$.
Cela peut aussi être vu comme un test rudimentaire : \og{}Si $I=\lone$ faire ceci, sinon faire cela.\fg{}

\myfigure{1}{
	\tikzinput{portes_10}
}
\end{exemple}

%--------------------------------------------------------------------
\subsection{Arbres}
\index{arbre logique}

On peut aussi représenter les connecteurs par des \defi{arbres} :

\myfigure{1}{
	\tikzinput{arbres_01}
}


\begin{exemple}
Considérons la formule :
\[	\big( P \land (Q \lor \lnot R) \big) \limplies (\lnot P \land R) \]
L'arbre correspondant est :
\myfigure{1}{
	\tikzinput{arbres_02}
}	
\end{exemple}

L'intérêt est que la structure d'arbre est adaptée au traitement algorithmique.
Un parcours préfixe en profondeur considère la racine, puis explore le sous-arbre de gauche, puis le sous-arbre de droite.

\medskip

On peut aussi adopter la \defi{notation polonaise}\index{notation polonaise}, c'est-à-dire écrire $\land P Q$ (comme une fonction de deux variables $\land(P,Q)$) au lieu de $P \land Q$.
En notation polonaise, la formule de l'exemple précédent s'écrit :
\[
\limplies \land P \lor Q \lnot R \land \lnot P R 
\]
C'est une expression particulièrement adaptée au traitement par ordinateur, car elle ne nécessite pas de parenthèses.
Pour un humain, c'est plus lisible de l'écrire ainsi :
\[
\limplies \Big( \land \big(P\big) \big(\lor (Q) (\lnot R)\big) \Big)  \Big(\land (\lnot P) (R)\Big) 
\]


%--------------------------------------------------------------------
\subsection{Formules algébriques}
\index{formule!algebrique@algébrique}

On note $\Bb = \{ \lzero, \lone \}$ les valeurs booléennes. 
Dans cette section, on fait les calculs modulo $2$, c'est-à-dire qu'on écrit $\lone+\lone = \lzero$ (en toute rigueur on devrait écrire $\lone + \lone \equiv \lzero \pmod 2$).
Noter que modulo $2$, on a $+\lone = - \lone$.

On note $x,y \in \{ \lzero, \lone \}$.
Voici les expressions algébriques qui transposent l'action des connecteurs :
\[
\begin{array}{l}
\lnot x = \lone + x \\
x \land y = x \times y \\
x \lor y = x + y - xy \qquad \text{ ou inclusif}\\
x \oplus y = x + y \qquad \text{ ou exclusif }	
\end{array}
\]
Cela se justifie facilement, par exemple, si $x=\lzero$, alors $\lone+x=\lone$, et si $x=\lone$ alors $\lone+x=\lone+\lone=\lzero$. Donc $\lone+x$ correspond bien à $\lnot x$. 

De même $x \times y$ ne peut valoir $\lone$ que si $x=\lone$ et $y=\lone$.

\begin{exemple}
	Ces expressions algébriques peuvent fournir des preuves. 
	Par exemple la double négation se justifie par :
	\[
	\lnot \lnot x = \lnot(\lone + x) = \lone + (\lone +x) = \lone+\lone+x = x
	\]
	Voici une preuve de la loi de De Morgan sur la négation d'un \tand, $\lnot( x \land y ) = (\lnot x) \lor (\lnot y)$ :
	\[ (\lnot x) \lor (\lnot y)
	= (\lone+x) \lor (\lone +y) 
	= (\lone+x) + (\lone +y) - (\lone+x)(\lone +y)
	= x +y - \lone -x -y -xy
	= \lone + xy
	= \lnot(xy)
	\]
	On a utilisé $-\lone=+\lone$, donc $-xy=xy$.
\end{exemple}



%%%%%%%%%%%%%%%%%%%%%%%%%%%%%%%%%%%%%%%%%%%%%%%%%%%%%%%%%%%%%%%%%%%%%
\section*{Exercices}

\subsection*{Vrai/faux, connecteurs}

\begin{exercicecours}
	Soit $P$ : \og{}$2 \times 2 = 5$\fg{} et $Q$ : \og{}$4+4 > 7$\fg{}.
	Dire si les propositions suivantes sont vraies ou fausses :
	$\lnot P$, $\lnot Q$, $P \land Q$, $P \lor Q$, $P \limplies Q$, $Q \limplies P$, $P \lequiv Q$.
\end{exercicecours}

\begin{exercicecours}
	On suppose $P$ vraie. Pour quelles valeurs de $Q$ les propositions suivantes sont-elles vraies : $P \land Q$, $P \lor Q$, $P \limplies Q$, $Q \limplies P$, $P \lequiv Q$ ? Même question si cette fois on suppose $P$ fausse.
\end{exercicecours}


\subsection*{Formules}

\begin{exercicecours}
Parmi les mots suivants lesquels sont des formules (au sens strict de la définition) :
\[\lnot \lnot P\qquad (\lnot P \lor Q \land R) \qquad (P \limplies (Q  \land (R \land S))) \qquad (P, Q \limplies R) \qquad ( (P \land \lnot Q \limplies) (Q \lor R))
\]	
\end{exercicecours}

\begin{exercicecours}
	Dresser la table de vérité pour $\lnot(P \lor Q)$ et celle pour $(\lnot P) \land (\lnot Q)$. En déduire la loi de De Morgan pour $\lor$. 
\end{exercicecours}

\begin{exercicecours}
	Prouver les formules d'associativité et de distributivité suivantes :
	\[ P \lor (Q \lor R) \equiv (P \lor Q) \lor R \qquad \qquad
	P \land (Q \lor R) \equiv (P \land Q) \lor (P \land R) \] 
\end{exercicecours}

\begin{exercicecours}
	Montrer les équivalences logiques suivantes :
	\[
	P \limplies Q \land R \equiv (P \limplies Q) \land (P \limplies R) \qquad
	P \lor Q \limplies R \equiv (P \limplies R) \land (Q \limplies R) \qquad
	P \lequiv Q \equiv (P \lor Q) \limplies (P \land Q)
	\]
	Traduire ces équivalences par une phrase en français.
\end{exercicecours}


\subsection*{Connecteurs (suite)}

\begin{exercicecours}
	Montrer que les formules suivantes sont des tautologies (elles sont toujours vraies, quelles que soient les valeurs de $P$ et $Q$) :
	\[
	P \lor (P \limplies Q) \qquad
	(P \land Q) \limplies P \qquad
	(P \limplies Q) \lor (P \limplies \lnot Q) \qquad
	(P \limplies Q) \lor (\lnot P \limplies Q)
	\]
\end{exercicecours}

\begin{exercicecours}
	Pour chacune des formules suivantes, trouver une formule d'expression la plus simple possible qui lui est logiquement équivalente :
	\[ Q \lor (P \land Q) \qquad
	P \land (P \lor Q) \qquad
	P \land (\lnot P \lor Q) \qquad
	(P \lor Q) \limplies (P \land Q)
	\]
\end{exercicecours}

\begin{exercicecours}
	Soit $x \in \Rr$, $P$ : \og{}$x \ge 2$\fg{} et $Q$ : \og{}$x^2 \ge 4$\fg{}.
	Trouver toutes les valeurs $x \in \Rr$, telles que $P \limplies Q$ soit vraie. Même question avec $Q \limplies P$, puis $P \lequiv Q$.
\end{exercicecours}

\begin{exercicecours}
	~
	\begin{enumerate}
	\item Déterminer la table de vérité de $P \limplies (Q \limplies R)$, puis celle de $(P \limplies Q) \limplies R$. Conclusion ?
	
	\item Montrer que $P \lequiv (Q \lequiv R) \equiv (P \lequiv Q) \lequiv R$.
	
	\item Montrer que $P \limplies (Q \limplies R) \equiv (P \land Q) \limplies R \equiv Q \limplies (P \limplies R)$.
	\end{enumerate}
\end{exercicecours}

\begin{exercicecours}
	On rappelle que $P \oplus Q$ désigne le \og{}ou exclusif\fg{} (c'est vrai lorsque $P$ est vrai, ou $Q$ est vrai, mais pas les deux).
	\begin{enumerate}
	\item Calculer la table de vérité de $P \oplus Q$.
	
	\item Montrer que $P \oplus Q$ revient à \og{}$P \neq Q$\fg{}, c'est-à-dire qu'on a $P \oplus Q \equiv (P \land \lnot Q) \lor (\lnot P \land Q)$. Quel est le lien entre $\oplus$ et $\lequiv$ ?
	
	\item Montrer que $\oplus$ est commutative et associative.
	
	\item Exprimer $P \land (Q \oplus R)$ uniquement avec les connecteurs $\land$, $\lor$ et $\lnot$. Idem avec $P \oplus (Q \land R)$.
	\end{enumerate}
\end{exercicecours}

\begin{exercicecours}
	Soit $P_1,P_2,\ldots,P_n$, avec $n \ge 1$.
	Quand est-ce que $P_1 \oplus P_2 \oplus \cdots \oplus P_n$ est vraie ?
\end{exercicecours}


\subsection*{Formules (suite)}

\begin{exercicecours}
	Les formules suivantes sont-elles satisfaisables ? Sont-elles des tautologies ?
	\[
	P \land Q \limplies P \qquad
	(P \lor Q) \land (\lnot P \lor \lnot Q) \qquad
	P \lor (Q  \land \lnot Q) \lequiv P \qquad
	(P \limplies Q) \land (Q \limplies \lnot P)
	\]
	\[
	\big( (P \limplies Q) \land (R \limplies S) \big) \limplies \big( (P \lor R) \limplies (Q \lor S) \big)
	\]
\end{exercicecours}

\begin{exercicecours}
	Montrer que les formules suivantes sont des tautologies :
	\begin{enumerate}
		
		\item $P \land (P \limplies Q) \limplies Q$ (\emph{modus ponens})
		
		\item $(P \limplies Q) \land \lnot Q \limplies \lnot P$ (\emph{modus tollens})
		
		\item $(P \limplies Q) \lequiv (\lnot Q \limplies \lnot P)$ (contraposition)
		
		\item ($P \limplies Q) \lequiv \big( (P \land \lnot Q) \limplies \bot \big)$ (raisonnement par l'absurde)
		
		\item $\big((P \limplies Q) \land (Q \limplies R)\big) \limplies (P \limplies R)$ (syllogisme/transitivité de l'implication).
		
		\item $(P \land \top) \lequiv P$
	\end{enumerate}
\end{exercicecours}

\begin{exercicecours}
	~
	\begin{enumerate}
		\item Construire le circuit pour $P \land (Q \lor R)$ et celui pour $(P \land Q) \lor (P \land R)$. Vérifier qu'ils produisent les mêmes sorties.
		
		\item Montrer qu'en utilisant uniquement des portes \tconnector{nand}, on peut construire des circuits dont les sorties sont identiques à l'action de $\lnot$, $\land$, $\lor$.
	\end{enumerate}
\end{exercicecours}

\begin{exercicecours}
	Pour chacune des formules suivantes, construire l'arbre associé et donner l'expression selon la notation polonaise :
	\[
	\lnot P \land \big( Q \lor (P \land R)\big) \qquad\qquad
	\big( (P \lor Q) \land \lnot R \big) \limplies \big( P \land (\lnot Q \land R) \big)
	\]
\end{exercicecours}

\begin{exercicecours}
	~
	\begin{enumerate}
		\item Justifier la formule algébrique $x \lor y = x+y-xy$.
		
		\item À l'aide des formules algébriques, montrer :
		\[ x \land (y \land z) = (x \land y) \land z \qquad
		x \lor (y \lor z) = (x \lor y) \lor z\qquad
		x \land (y \lor z) = (x \land y) \lor (x \land z)
		\]
		
		\item Calculer $x \land (y \oplus z)$ et $(x \land y) \oplus (x \land z)$. Conclusion ?
		
		\item Trouver une formule algébrique pour $x \limplies y$. Idem pour $x \lequiv y$.
		
		\item À l'aide des formules algébriques, montrer $x \limplies y = \lnot y \limplies \lnot x$.
	\end{enumerate}	
\end{exercicecours}


\end{document}
